\section{Introduction}

Light field (LF) imaging captures both spatial and angular information of light rays, enabling advanced 3D vision tasks. Depth estimation from LF images is a core component, facilitating applications in autonomous navigation, virtual reality, and robotic manipulation [1]. By exploiting the epipolar plane image (EPI) (EPI) structure, where a 3D point projects as a straight line across angular views, existing methods achieve high accuracy in ideal environments. 

However, obtaining accurate depth for non-Lambertian surfaces, such as specular, translucent, or scattering materials, remains a challenging problem for both traditional algorithms and deep networks [2]. The primary difficulty lies in the violation of the Lambertian assumption. Under standard physical models, pixel intensity remains consistent across viewpoints, forming clear linear structures in EPIs. In contrast, non-Lambertian reflections shift dynamically with the viewpoint, destroying the EPI linearity and introducing severe structural ambiguities, such as X-shaped patterns or blurred slopes [3]. 

Despite substantial efforts, most current LF depth estimation methods primarily support Lambertian scenes, making them unreliable for mixed or highly reflective environments [4]. When applied to non-Lambertian regions, the performance of state-of-the-art EPI-based networks degrades significantly. For instance, the mean absolute error (MAE) in non-Lambertian regions often exceeds 0.40, significantly surpassing the acceptable threshold for practical deployment. Furthermore, addressing this issue is severely hindered by data scarcity. Public non-Lambertian LF datasets contain only a handful of training scenes (e.g., merely 4 scenes in specific benchmarks), creating a formidable data barrier for purely data-driven deep learning models [5]. Additionally, complex textures in standard Lambertian scenes cause EPI slope ambiguity, exposing the inherent architectural limitations of relying solely on EPI features.

To address these physical and data-driven bottlenecks, we propose a Unified Dual-Mask Physical Model for non-Lambertian LF depth estimation. Instead of relying on conventional 2D discrete Fourier transform (DFT) which struggles under low angular resolutions, we introduce a three-layer angular signal decomposition mechanism. This model decouples the angular signal into DC (ambient illumination), linear (parallax/depth), and residual (material/reflectance) components, enabling robust pixel-level material classification via geometric angular features. Building upon this physical insight, we design a GeometricDualMask network that adaptively routes LF features into specialized branches for Lambertian and non-Lambertian regions. To overcome the severe data scarcity, we incorporate a domain-balanced sampling strategy with a weighted random sampler, ensuring stable multi-domain joint training.

Our main contributions are summarized as follows:
1) We propose a three-layer angular signal decomposition model and a geometric feature-based material classification mechanism. This explicitly reveals the physical causes of EPI linearity violation and provides a robust material identification scheme under low angular resolutions.
2) We develop a GeometricDualMask network architecture that utilizes dual masks to route features into domain-specific branches. This effectively mitigates depth estimation collapse caused by non-Lambertian EPI structure destruction.
3) We introduce a domain-balanced sampling strategy to alleviate the overfitting issue caused by extreme non-Lambertian data scarcity. This achieves superior overall and mixed-scene depth estimation accuracy compared to existing baselines.

The remainder of this paper is organized as follows. Section II reviews related work on LF depth estimation. Section III details the proposed physical model and network architecture. Section IV presents the experimental setup and comprehensive results. Finally, Section V concludes the paper.

Traditional light field depth estimation algorithms primarily rely on strict physical assumptions, such as Lambertian reflectance and photo-consistency [1]. Methods based on epipolar plane image (EPI) (EPI) slopes or structure tensors evaluate local linear structures to infer depth [6]. However, these techniques fail on non-Lambertian surfaces because specular or translucent reflections dynamically shift across viewpoints. This shifting destroys EPI linearity and creates ambiguous X-shaped patterns, leading to severe depth errors [7]. Furthermore, correspondence-based or defocus-based approaches struggle in textureless or highly reflective regions. They often produce incorrect matches and suffer from high computational complexity, limiting their practical applicability in mixed-material environments [8].

To overcome the limitations of handcrafted features, early deep learning methods formulate depth estimation as an end-to-end regression task [9]. Networks utilizing 3D convolutional neural networks or basic EPI architectures extract hierarchical features directly from raw light field data [10]. Despite their success in ideal scenes, these models implicitly assume Lambertian properties during training. Consequently, they lack explicit material awareness and treat view-dependent intensity variations as valid geometric cues [11]. When applied to non-Lambertian regions, these networks mistakenly interpret specular highlights or refractions as depth discontinuities. This physical misinterpretation causes significant performance degradation, with mean absolute errors often exceeding acceptable thresholds in complex scenes [12].

Recent advanced approaches attempt to address complex reflections by introducing attention mechanisms, dual-branch architectures, or auxiliary defocus heads [13]. Some methods employ frequency-domain analyses, such as 2D discrete Fourier transform, to separate material properties from geometric structures . However, these frequency-based techniques suffer from severe spectral aliasing under low angular resolutions. In practical light field cameras with sparse angular sampling, the frequency features overlap, rendering material classification highly unreliable [14]. Moreover, existing multi-domain or separate per-domain models fail to resolve the inherent data scarcity problem. Public non-Lambertian datasets contain extremely limited training scenes, causing complex networks to overfit rapidly and preventing unified inference across mixed environments [15].

In summary, current methods either ignore the physical degradation of EPI structures or rely on frequency features that fail at practical angular resolutions. Furthermore, they lack a unified mechanism to handle mixed materials without overfitting to scarce non-Lambertian data. Therefore, there is an urgent need for a unified approach that explicitly decouples material properties from geometric cues and adaptively routes features for robust depth estimation in mixed scenes.

In this paper, we propose a unified dual-mask physical model, named GeometricDualMask, for robust light field depth estimation in mixed-material environments. Instead of relying on implicit Lambertian assumptions, our approach explicitly decouples angular signals into illumination, disparity, and material residual components. This physical decomposition enables precise pixel-level material classification, which subsequently routes the extracted light field features into two specialized network branches. By separately processing Lambertian linear structures and non-Lambertian X-shaped patterns, the proposed framework effectively mitigates depth estimation collapse caused by specular reflections.

The main contributions of this paper are summarized as follows:

\begin{itemize}
 \item \textbf{Contribution 1}: We propose a three-layer angular signal decomposition mechanism integrated with geometric material classification. By decoupling angular signals into illumination, linear, and residual components, we extract geometric features to train a random forest classifier. This overcomes the frequency aliasing of 2D-DFT at low angular resolutions, yielding robust and interpretable pixel-level material masks.
 \item \textbf{Contribution 2}: We develop the GeometricDualMask network architecture to explicitly handle complex mixed reflection scenarios. This dual-branch model utilizes the generated masks to route features into a Lambertian branch with continuous depth priors and a non-Lambertian branch with cross-guided geometric matching. This design resolves the failure of single EPI assumptions on glossy surfaces, significantly enhancing depth accuracy across complex scenes.
 \item \textbf{Contribution 3}: We introduce a domain-balanced training strategy integrated with extensive evaluations on mixed-material scenes. By employing weighted random sampling and tone-mapping augmentation, we effectively alleviate the severe overfitting caused by scarce non-Lambertian training data. This ensures high generalization across diverse datasets, establishing a comprehensive and reliable benchmark for mixed-material light field depth estimation.
\end{itemize}
Extensive experiments on the HCInew, UrbanLF-Syn, and Non-Lambertian datasets clearly demonstrate the superiority of our proposed method. Utilizing 4-direction EPI inputs and a learning rate of $1 \times 10^{-4}$, our model achieves an overall Mean Absolute Error (MAE) of 0.133 and a mixed urban MAE of 0.081. These results substantially outperform existing baselines, confirming the effectiveness of our physical decomposition and dual-mask routing strategy in practical applications. To contextualize these advancements, we first survey the traditional light field depth estimation methods that have historically relied on similar physical cues.